% Thesis Abstract -----------------------------------------------------
% BMRC style

%\begin{abstractslong}    %uncommenting this line, gives a different abstract heading
\begin{abstracts}        %this creates the heading for the abstract page.

\textbf{Background:} Chronic disease and functional decline are major health challenges for older adults, yet how well self-reported health burden can predict disease onset across longitudinal waves remains underexplored.

\textbf{Methods:} A longitudinal machine learning framework was developed using five waves of TILDA data (2009--2019). Four model families, \texttt{logistic Regression}, \texttt{XGBoost}, \texttt{random Forest}, and a \texttt{multi-Layer Perceptron}, were compared across a single wave baseline and two pooled multi wave designs. Performance was assessed using \gls{CV} \gls{AUC}. The contribution of individual domain-specific predictors was evaluated using a bootstrapped add-on analysis.

\textbf{Results:} Pooled multi wave modelling improved predictive performance and stability over the single wave baseline, increasing \gls{CV} \gls{AUC} from $0.589 \pm 0.072$ to $0.718 \pm 0.006$. Reducing the feature set to the 20 most important predictors caused no measurable performance loss, while cognitive, psychosocial, nutritional, and behavioural add-on variables produced no statistically significant improvement.

\textbf{Conclusions:} Among community dwelling older adults, chronic disease risk prediction was driven primarily by accumulated self-reported and functional health measures, particularly medication use, pain, and self-rated health, rather than domain specific biomarker or lifestyle profiles. These findings suggest that publicly available survey data can support meaningful risk stratification with limited feature sets.

\textbf{Keywords:} ageing, \glsentryfull{ML}, predictive modelling, chronic disease, functional decline, health burden, \glsentryfull{TILDA}, \glsentryfull{XGB}

\end{abstracts}
%\end{abstractlongs}


% ---------------------------------------------------------------------------
% ----------------------- end of thesis sub-document ------------------------
% ---------------------------------------------------------------------------